
\paragraph{Uncertainty Quantification for Deep Time-Series Forecasting} The two major paradigms of  Uncertainty Quantification (UQ) methods for deep neural networks are Bayesian and Frequentist. Bayesian approaches estimate a distribution over the model parameters given data, and then marginalize these parameters to form output distributions via Markov Chain Monte Carlo (MCMC) sampling \cite{welling2011bayesian, neal2012bayesian, chen2014stochastic} or variational inference (VI)  \cite{graves2011practical,kingma2015variational, blundell2015weight, louizos2017multiplicative}. \cite{wang2019deep, wu2021quantifying} propose Bayesian Neural Network (BNN)-based models for UQ of spatiotemporal forecasts. In practice, Bayesian UQ has two major drawbacks: (1) it requires significant modification to the training process of deep models, and (2) BNNs can be computationally expensive and difficult to optimize, especially for larger networks \cite{lakshminarayanan2017simple, zadrozny2001obtaining}. UQ for deep neural network time series forecasts often adopt approximate Bayesian inference such as MC-dropout \cite{gal2016theoretically, gal2017concrete}.

Frequentist UQ methods emphasize robustness against variations in the data. These approaches either rely on resampling the data or learning an interval bound to encompass the dataset. For time series forecasting UQ, frequentist approaches include ensemble methods such as bootstrap \cite{efron2016, alaa2020frequentist} and jackknife methods  \cite{jackknife+after, bjrnn}; interval prediction methods include interval regression through proper scoring rules \cite{kivaranovic2020adaptive, wu2021quantifying}, \edit{and quantile regression \cite{ takeuchi2006nonparametric}, with many recent advances for time series UQ  \cite{tagasovska2019single, gasthaus2019probabilistic,park2022learning, kan2022multivariate}.} Many of the frequentist methods produces asymptotically valid confidence regions and can be categorized as distribution-free UQ techniques as they are (1) agnostic to the underlying model and (2) agnostic to data distribution.

\paragraph{Conformal Prediction.} 
%, and has shown promise in many realistic applications \cite{eklund2015application, angelopoulos2022image, lei2021conformal}
Conformal prediction (CP) is an important member of distribution-free UQ methods; we refer readers to \cite{angelopoulos2021gentle} for a comprehensive introduction and survey. CP has become popular because of its simplicity, theoretical soundness, and low computational cost. A key difference between CP and other UQ methods is that under the \edit{exchangeability} assumption, conformal methods guarantee validity in finite samples \cite{vovk2005algorithmic}.

Most relevant to our work is recent progress on expanding conformal prediction to time-series forecasting. According to \cite{Schaar2021conformaltime} there are two settings: when the data is generated from one single time series or from multiple independent time series. For the first setting, ACI \cite{gibbs2021adaptive} and EnbPI \cite{xu2021conformal} developed CP algorithms that relaxes the \edit{exchangeability} assumption while maintaining asymptotic validity via online learning (former) and ensembling (later); \cite{zaffran2022adaptive} further improves on the online adaptation, and \cite{sousa2022general} combines EnbPI with conformal quantile regression \edit{\cite{romano2019conformalized}} to model heteroscedastic time series. These algorithms for single time series are not designed to model multiple independent time series where extracting common patterns can improve forecast and UQ results. The validity guaranteed by these methods are on \textit{average} over \edit{time steps} and are often asymptotic, rather than covering the full horizon as in our setting.  
\cite{Schaar2021conformaltime} shares with us the multi-step forecasting for multiple time \edit{series} setting, though they focuses on univariate medical time series. We show that their method of applying Bonferroni correction produces inefficient confidence regions, especially when data is multidimensional. \cite{neeven2018weather} and \cite{messoudi2020multitarget} are CP algorithms for multi-target regression in the non-temporal setting, creating box-like regions to account for the correlations between the labels. In our work we develop CoupulaCPTS for multivariate multi-step time series forecasting.

